\newpage
\section*{Заключение}
\addcontentsline{toc}{section}{Заключение}

В ходе выполнения практической работы были рассмотрены основные
принципы адресации в компьютерных сетях, расчета параметров
беспроводных соединений Wi-Fi и определения времени передачи данных
по сети.

В первом задании исходная сеть \texttt{168.60.20.0/22}, содержащая
1024 адреса, была разделена на 5 подсетей. Были получены две подсети
по 128 адресов с маской \texttt{/25} и три подсети по 256 адресов с
маской \texttt{/24}. Таким образом, все адресное пространство
исходной сети было распределено между пятью подсетями, а требование
о наличии как минимум двух подсетей по 128 адресов было выполнено.

Во втором задании были определены параметры Wi-Fi соединений двух
устройств стандарта 802.11ac. Для устройства A был выбран MCS 8 с
модуляцией 256-QAM и скоростью 324 Мбит/с при GI = 800 нс и
360 Мбит/с при GI = 400 нс. Для устройства B был выбран MCS 3 с
модуляцией 16-QAM и скоростью 216 Мбит/с при GI = 800 нс и
240 Мбит/с при GI = 400 нс. Несмотря на большее количество
пространственных потоков MIMO у устройства B, его скорость соединения
ниже из-за более низкого уровня сигнала и использования менее
эффективной модуляции. В результате скорость устройства A в 1,5 раза
выше скорости устройства B.

В третьем задании было рассчитано время передачи файла размером
83 Мбайт по сети с использованием Ethernet DIX и IPv4. С учетом
заголовков и контрольной суммы FCS полезная нагрузка одного кадра
составила 924 байта, для передачи файла потребовалось 94\,191 кадр.
С учетом заданной доли повторных передач было выполнено 730 повторных
передач. Итоговое время передачи файла составило приблизительно
8\,666{,}969 с, или 2 ч 24 мин 27 с.

Таким образом, в ходе лабораторной работы были закреплены навыки
расчета адресного пространства и разбиения IP-сетей на подсети,
определения параметров и скорости Wi-Fi соединений, а также расчета
времени передачи данных с учетом сетевых заголовков и повторных
передач пакетов.